\documentclass[11pt]{article}
\usepackage[letterpaper,margin=1in]{geometry}
\usepackage[T1]{fontenc}
\usepackage{lmodern,microtype}
\usepackage{amsmath,amssymb,amsthm,mathtools,mathrsfs}
\usepackage{enumitem,booktabs,array,fancyhdr,placeins}
\usepackage{xcolor}
\usepackage[colorlinks=true,linkcolor=blue!40!black,citecolor=blue!40!black,urlcolor=blue!40!black]{hyperref}
\usepackage{bookmark}
\hypersetup{pdftitle={Memory-constrained minimax first-order optimization},
 pdfsubject={An exact semialgebraic decision procedure and two convex benchmarks}}
\newcommand{\supportlabel}[1]{\textcolor{blue!60!black}{\small\textsf{[#1]}}}
\newcommand{\eventproof}{\supportlabel{Event proof}}
\newcommand{\literature}{\supportlabel{Literature}}
\newcommand{\computation}{\supportlabel{Computation}}
\pagestyle{fancy}
\fancyhf{}
\fancyhead[L]{\small Memory-constrained minimax optimization}
\fancyhead[R]{\small Finite query horizons}
\fancyfoot[C]{\thepage}
\setlength{\headheight}{14pt}
\setlength{\parskip}{1.5pt}
\setlength{\emergencystretch}{2em}
\numberwithin{equation}{section}
\newtheorem{theorem}{Theorem}[section]
\newtheorem{proposition}[theorem]{Proposition}
\newtheorem{lemma}[theorem]{Lemma}
\newtheorem{corollary}[theorem]{Corollary}
\theoremstyle{definition}
\newtheorem{definition}[theorem]{Definition}
\newtheorem{assumption}[theorem]{Assumption}
\newcommand{\R}{\mathbb R}
\newcommand{\K}{\mathbb R_{\mathrm{alg}}}
\newcommand{\W}{\mathcal W_N}
\newcommand{\A}{\mathcal A}
\newcommand{\Q}{\mathcal Q}
\newcommand{\risk}{\mathcal R}
\DeclareMathOperator{\rank}{rank}
\DeclareMathOperator{\im}{im}
\DeclareMathOperator{\diag}{diag}
\DeclareMathOperator{\Sym}{Sym}
\begin{document}
\begin{center}
{\LARGE\bfseries Memory-constrained minimax\\[3pt] first-order optimization\par}
\vspace{0.65em}
{\large An exact semialgebraic decision procedure\\
and two convex benchmarks\par}
\vspace{0.7em}
Francisco Hern\'andez Iglesias and Jack Mayo\\[2pt]
{\small \textit{kurtos.ai} --- Computation/Applied Math}\\[4pt]
{\small Repository: \url{https://github.com/kurtos-ai/mathathon32-memory} (private)}\\[2pt]
{\small Lean formalization commit: \texttt{f5aeb9cdfd0a777dc34480dd187c994f7b13b737}}\\[2pt]
{\small Event work: Nov~13, 8PM -- Nov~15, noon Pacific, 2026}\\[2pt]
November 2026
\end{center}
\thispagestyle{plain}
\begin{abstract}
We consider deterministic first-order methods with predetermined scalar coefficients, a prescribed finite number of exact gradient queries, and a bounded number of persistent vector registers. We prove an exact characterization of memory by ranks of past-to-future coefficient blocks. Under an effective finite semialgebraic description of attainable losses, we prove that projecting the performance condition, imposing the memory equations, and deciding their intersection computes the minimax values and all strict memory gains. Every finite value is a computable real algebraic number. We extend the rank formula and decision procedure to subgradient methods with a projection after every increment, under an explicitly specified register model. For convex quadratics, we calculate the exact value and show that one vector register suffices for every horizon. For convex smooth functions, we give the finite interpolation formula, recover the optimal two-register guarantee, and determine the strict first-register separation using a cited uniqueness theorem and a compactness argument.
\end{abstract}

\noindent\textit{Literature note.} Two works named by the problem frame our study:
d'Aspremont--Scieur--Taylor \cite{AST}, \emph{Acceleration Methods} (2021), Ch.~4 (performance estimation and semialgebraic formulations) and
Das Gupta--Van Parys--Ryu \cite{BnBPEP}, BnB-PEP (2024) (branch-and-bound for exact worst-case analysis).
Our results build on the interpolation framework of Taylor--Hendrickx--Glineur \cite{THG}
and the OGM optimality theory of Kim--Fessler \cite{KF} and Drori \cite{Drori},
with the uniqueness theorem of Grimmer--Jo--Park \cite{GJP} for the smooth-convex separation.
All other claims below are proved here \eventproof\ unless marked \literature\ or \computation.

\section{Problem statement and assumptions}
\label{sec:problem}
\subsection{Instances, oracle, and coefficient maps}
Fix a known integer $N\ge1$, an initial radius $R>0$, and a family of instance sets with nonempty union
\[
 \mathscr I=\bigcup_{n\ge1}\mathscr I_n.
\]
An instance in $\mathscr I_n$ consists of a continuously differentiable function $f:\R^n\to\R$, a starting point $x_0$, and a minimizer $x_\star$ satisfying $\|x_0-x_\star\|\le R$. Additional restrictions on $f$ define the chosen function class. The algorithm receives $x_0$ and access to the exact oracle $y\mapsto\nabla f(y)$. The minimizer is used only to specify the class and the loss; it is not supplied to the algorithm. The loss is
\[
 E(W,I)=f(y_N)-f(x_\star)\ge0.
\]
The domain is always $\R^n$. All arithmetic is exact.

\begin{definition}[Methods and query budget]\label{def:methods}
Let $\W$ be the vector space of strictly lower triangular real $(N+1)\times(N+1)$ matrices, with rows and columns indexed by $0,\ldots,N$. Write $a_{ij}=W_{ij}$ for $j<i$; all entries with $j\ge i$ are zero. Thus $\W\cong\R^{N(N+1)/2}$. The method prescribed by $W$ is
\begin{equation}
 y_0=x_0,\qquad g_i=\nabla f(y_i)\quad(0\le i<N),\qquad
 y_i=x_0+\sum_{j<i}a_{ij}g_j\quad(1\le i\le N).
 \label{eq:history}
\end{equation}
The queries are $y_0,\ldots,y_{N-1}$ and the output is $y_N$. Thus there are exactly $N$ gradient queries. In particular, the intermediate points must be generated in order: $g_j$ becomes available only after querying at $y_j$.

All $a_{ij}\in\R$ are chosen before execution. They may depend on $N$ and the known class parameters, including $R$, but not on the dimension, the instance, or any oracle observations. No sign or magnitude restrictions are imposed.
\end{definition}

The same schedule is used in every finite dimension. Our method class is restricted to these predetermined linear combinations; it does not include adaptive scalar coefficients, nonlinear state encodings, or randomization.

\subsection{The memory model}
\begin{definition}[Realizations and register count]\label{def:memory}
An $s$-register realization of a fixed $W$, with $s\ge1$, is a tuple of predetermined constants
\[
 F_k\in\R^{s\times s},\quad b_k,c_k\in\R^s\ (0\le k<N),
 \qquad o\in\R^s
\]
such that the recurrence
\begin{align}
 Z_0&=e_1\otimes x_0,&
 \widetilde y_k&=(c_k^\top\otimes I_n)Z_k,\nonumber\\
 Z_{k+1}&=(F_k\otimes I_n)Z_k+(b_k\otimes I_n)g_k,&
 \widetilde y_N&=(o^\top\otimes I_n)Z_N
 \label{eq:state}
\end{align}
reproduces all the $y_i$ in~\eqref{eq:history} for every $n$, every $x_0\in\R^n$, and every sequence of independent formal inputs $g_0,\ldots,g_{N-1}\in\R^n$. The constants have the same instance-independence restrictions as the $a_{ij}$.

The state $Z_k\in(\R^n)^s$ is the entire persistent instance-dependent state. A fresh gradient and temporary work within an update are not counted as persistent registers. There is no free persistent copy of $x_0$, no stored transcript outside $Z_k$, and no additional data-dependent scalar state.
\end{definition}

Equality for formal inputs implies equality on oracle trajectories by induction. Define
\begin{equation}
 s(W)=\min\{s\ge1:W\text{ has an }s\text{-register realization}\},
 \qquad
 \A_m=\{W:s(W)\le m+1\}\quad(m\ge0).
 \label{eq:memory-class}
\end{equation}
The minimum is over implementations of the fixed array $W$. It exists: one can store $x_0$ and every past gradient. We use $m$ for the number of auxiliary registers, so $m=0$ permits one persistent vector. The time index and predetermined scalar constants are available without being counted as vector registers.

\subsection{Minimax values and the finite semialgebraic input}
\begin{definition}[Risk and memory improvement]\label{def:values}
Set
\begin{equation}
 \risk(W)=\sup_{I\in\mathscr I}E(W,I),\qquad
 V_m=\inf_{W\in\A_m}\risk(W),\qquad
 V_{\mathrm{lin}}=\inf_{W\in\W}\risk(W).
 \label{eq:values}
\end{equation}
The supremum includes all finite dimensions. These quantities belong to $[0,+\infty]$. An additional auxiliary register helps at budget $m$ if $V_{m+1}<V_m$. When $V_{\mathrm{lin}}<\infty$, the least sufficient auxiliary budget is
\begin{equation}
 m_\star=\min\{m\ge0:V_m=V_{\mathrm{lin}}\}.
 \label{eq:mstar}
\end{equation}
\end{definition}

The standing hypothesis is that, for a fixed query budget, all possible terminal losses can be described exactly using finitely many real variables and polynomial constraints. This gives a finite problem to optimize; it does not supply its optimal value. For effective computation, write $\K=\overline{\mathbb Q}\cap\R$ for the real algebraic numbers: real roots of nonzero polynomials with rational coefficients, represented exactly.

\begin{assumption}[Exact finite semialgebraic loss model]\label{ass:finite}
Fix the instance class and horizon $N$. The input supplies a finite integer $q$ and the following two explicit formulas in the method coefficients $W$ and auxiliary real variables $z=(z_1,\ldots,z_q)$:
\begin{enumerate}[label=(\roman*),leftmargin=2.3em,itemsep=2pt]
\item A feasibility condition $\Psi_N(W,z)$, formed from finitely many polynomial equalities and inequalities using \emph{and}, \emph{or}, and \emph{not}. A vector $z$ is called feasible for $W$ when this condition holds.
\item A polynomial expression $e_N(W,z)$ giving the loss represented by a feasible vector $z$.
\end{enumerate}
All fixed coefficients in these formulas are known elements of $\K$. The integer $q$ and the formulas may depend on $N$ and the prescribed class, but they are the same for every real method $W$ and every instance dimension $n$.

Exactness requires both of the following statements for every fixed $W$:
\begin{enumerate}[label=(\alph*),leftmargin=2.3em,itemsep=2pt]
\item \emph{No actual loss is omitted.} For every admissible instance $I\in\mathscr I$, there is a feasible $z\in\R^q$ such that $e_N(W,z)=E(W,I)$.
\item \emph{No fictitious loss is introduced.} For every feasible $z\in\R^q$, there is an admissible instance $I\in\mathscr I_n$ in some finite dimension $n$ such that $E(W,I)=e_N(W,z)$.
\end{enumerate}
Equivalently, the two sets of real loss values are equal:
\begin{equation}
 \{e_N(W,z):\Psi_N(W,z)\}
 =
 \{E(W,I):I\in\mathscr I\}.
 \label{eq:loss-model}
\end{equation}
For a family of input horizons, these data are supplied for each $N$ or are effectively constructible from $N$, without using unknown minimax values. Every fixed numerical class parameter occurring in the formulas is algebraic. The variables $W$ and $z$ themselves range over all real numbers.
\end{assumption}

The vector $z$ is used to describe possible losses in the analysis. It is not an algorithmic state, is not supplied to the method, and is not counted as memory; $q$ has no register-count interpretation. Nor must $z$ encode an entire function or identify a unique instance. Different instances may give the same represented loss. In the smooth-convex construction of Proposition~\ref{prop:smooth-model}, $z$ consists of a Gram matrix and function values.

For a fixed $W$, Assumption~\ref{ass:finite} gives precisely
\[
 \risk(W)=\sup_{\{z\in\R^q:\Psi_N(W,z)\}}e_N(W,z).
\]
The supremum on the right remains to be computed, as does the subsequent minimization over eligible methods $W$.

This assumption is stronger than merely describing the function class by polynomial conditions: the exact finite reduction, valid across all instance dimensions, must be supplied or proved. The objective functions themselves need not be semialgebraic. Propositions~\ref{prop:quadratic-risk} and~\ref{prop:smooth-model} verify the assumption for the two classes treated below. A semialgebraic loss expression can also be accommodated by adding its value as a variable and its graph to $\Psi_N$.

The computational problem is to determine $V_m$, all strict comparisons $V_{m+1}<V_m$, and $m_\star$, for the given finite horizon and input description.

\section{Memory as a finite system of rank equations}
\label{sec:memory}
For $0\le k\le N$, define the augmented cut matrix
\begin{equation}
 H_k(W)=
 \begin{pmatrix}
 1&a_{k0}&\cdots&a_{k,k-1}\\
 1&a_{k+1,0}&\cdots&a_{k+1,k-1}\\
 \vdots&\vdots&&\vdots\\
 1&a_{N0}&\cdots&a_{N,k-1}
 \end{pmatrix}.
 \label{eq:cut}
\end{equation}
Its columns index the past inputs $x_0,g_0,\ldots,g_{k-1}$; its rows index the remaining queries and output. At $k=0$ it is the column of ones.

\begin{theorem}[Exact realization size \eventproof]\label{thm:realization}
Every coefficient map satisfies
\begin{equation}
 s(W)=\max_{0\le k\le N}\rank H_k(W)
 \le 1+\lfloor N/2\rfloor.
 \label{eq:cut-rank}
\end{equation}
In particular $\A_{\lfloor N/2\rfloor}=\W$.
\end{theorem}
\begin{proof}
To prove $s(W)\ge\max_k\rank H_k(W)$, work in one scalar coordinate. The rank formula is a specialization of the time-varying realization theory of van der Veen \cite{vdVeen95} and Dewilde--van der Veen \cite{DewildevdVeen98}, here specialized to the Hankel-like cut matrices $H_k$; we give the elementary proof. Before query $k$, any $s$-register state is a linear image of $(x_0,g_0,\ldots,g_{k-1})^\top$. With future inputs set to zero, all remaining outputs are linear observations of that state. Hence $H_k$ factors through $\R^s$, and $\rank H_k\le s$.

To prove $s(W)\le\max_k\rank H_k(W)$, let $E_k=\im H_k\subseteq\R^{N-k+1}$. A canonical state is the vector of all past contributions to remaining outputs. Its first coordinate is the current query. Let $D_k$ delete that coordinate and put
\[
 h_k=(a_{k+1,k},\ldots,a_{N,k})^\top.
\]
The identity $H_{k+1}=[\,D_kH_k\ \ h_k\,]$ implies
$D_kE_k\subseteq E_{k+1}$ and $h_k\in E_{k+1}$. Thus the update
\[
 v_{k+1}=D_kv_k+h_kg_k
\]
maps the canonical state into $E_{k+1}$ and generates the next required past contributions. At the last cut its only coordinate is the terminal output.

Let $s=\max_k\dim E_k$. Choose linear injections $J_k:E_k\to\R^s$ and left inverses $L_k:\R^s\to E_k$. Set $F_k=J_{k+1}D_kL_k$, $b_k=J_{k+1}h_k$, and use the first coordinate of $L_k$ as the observation. Choose $J_0$ to send the all-ones vector spanning $E_0$ to $e_1$. These maps realize $W$ with the required initialization. They depend only on $W$, and tensoring with $I_n$ proves the claim in every dimension.

Finally, $H_k$ has size $(N-k+1)\times(k+1)$, so its rank is at most $\min(N-k+1,k+1)$. The maximum of this bound is $1+\lfloor N/2\rfloor$.
\end{proof}

\begin{proposition}[Normalized-register equivalence \eventproof]\label{prop:normalized}
Fix $W$ and an integer $s\ge1$. An $s$-register realization exists if and only if
$W$ admits a realization with the same number of registers satisfying
\[
 Z_0=\mathbf1_s\otimes x_0,\qquad
 y_k=(e_1^\top\otimes I_n)Z_k\quad(0\le k\le N),\qquad
 F_k\mathbf1_s=\mathbf1_s\quad(0\le k<N).
\]
Thus initializing every register at $x_0$, using the first register as the query
and terminal output, and requiring transition rows to sum to one do not change
the minimum register count $s(W)$.
\end{proposition}
\begin{proof}
Suppose an $s$-register realization exists. By Theorem~\ref{thm:realization},
$E_k=\im H_k(W)$ has dimension at most $s$. Let $q_k$ be its all-ones vector
and $\ell_k$ its first-coordinate functional; then $\ell_k(q_k)=1$.
Choose a basis of $E_k$ consisting of $q_k$ and a basis of $\ker\ell_k$.
Map $q_k$ to $\mathbf1_s$ and the remaining basis vectors to linearly independent
vectors in $\ker(e_1^\top)\subset\R^s$. This defines an injection $J_k$ with
$J_kq_k=\mathbf1_s$ and $e_1^\top J_k=\ell_k$.

Extend the chosen vectors to a basis of $\ker(e_1^\top)$. The additional
vectors span a complement $C_k$ of $J_k(E_k)$ contained in that kernel.
Define $L_k$ to invert $J_k$ on $J_k(E_k)$ and to vanish on $C_k$. Then
\[
 L_kJ_k=\operatorname{id}_{E_k},\qquad
 L_k\mathbf1_s=q_k,\qquad \ell_kL_k=e_1^\top.
\]
With $D_k$ and $h_k$ from the canonical construction in
Theorem~\ref{thm:realization}, set $F_k=J_{k+1}D_kL_k$ and $b_k=J_{k+1}h_k$.
Since $D_kq_k=q_{k+1}$, we have $F_k\mathbf1_s=\mathbf1_s$.
Initialize the state at $J_0q_0\otimes x_0=\mathbf1_s\otimes x_0$.
The identity $L_kJ_k=\operatorname{id}$ shows inductively that these states
encode the canonical states. The identity $e_1^\top J_k=\ell_k$ makes their
first registers exactly the required queries and terminal output.

Conversely, choose an invertible $A\in\R^{s\times s}$ with
$A\mathbf1_s=e_1$. Conjugating a normalized state by $A\otimes I_n$ gives
initialization $e_1\otimes x_0$, transitions $AF_kA^{-1}$, input vectors
$Ab_k$, and readouts $e_1^\top A^{-1}$. This is an $s$-register realization
in Definition~\ref{def:memory} with the same outputs.
\end{proof}

We next separate coefficient directions that affect memory from directions along which memory is constant. Let $\mathfrak T_N$ be the space of strictly lower triangular $N\times N$ real matrices, with indices $0,\ldots,N-1$. Define
\begin{equation}
 d_j=a_{j+1,j},\qquad
 T_{ij}=a_{i+1,j}-a_{ij}\quad(j<i),\qquad T_{ij}=0\quad(j\ge i).
 \label{eq:coordinates}
\end{equation}
The inverse transformation is
\begin{equation}
 a_{ij}=
 \begin{cases}
 d_j+\displaystyle\sum_{\ell=j+1}^{i-1}T_{\ell j},
       &0\le j<i\le N,\\[3pt]
 0,    &0\le i\le j\le N.
 \end{cases}
 \label{eq:inverse}
\end{equation}
Thus the sum formula applies only below the diagonal. An empty sum is zero, so $a_{j+1,j}=d_j$.
Write the resulting map as $W(d,T)$. Thus $\W\cong\R^N\times\mathfrak T_N$. Equivalently, the additive group $\R^N$ acts by $a_{ij}\mapsto a_{ij}+h_j$; its orbits are the fibers with fixed $T$.
Here $W$ is already strictly lower triangular of size $(N+1)\times(N+1)$. The matrix $T$ is a different strictly lower triangular matrix of size $N\times N$, recording consecutive-row differences. The vector $d$ records the first subdiagonal entries $W_{j+1,j}$, not diagonal entries of $W$.

For $1\le k<N$, define the bottom-left block of $T$ by retaining rows $k,\ldots,N-1$ and columns $0,\ldots,k-1$:
\[
 S_k(T)=T_{\{k,\ldots,N-1\},\,\{0,\ldots,k-1\}}
 =
 \begin{pmatrix}
 T_{k,0}&\cdots&T_{k,k-1}\\
 T_{k+1,0}&\cdots&T_{k+1,k-1}\\
 \vdots&&\vdots\\
 T_{N-1,0}&\cdots&T_{N-1,k-1}
 \end{pmatrix}
 \in\R^{(N-k)\times k}.
\]
Each column corresponds to an already-observed gradient $g_j$ with $j<k$, and each row corresponds to a future increment $y_{i+1}-y_i$ with $i\ge k$. The entry $T_{ij}=a_{i+1,j}-a_{ij}$ is the coefficient of that past gradient in that increment. At $k=0$ or $k=N$, the block is empty and has rank zero.

\begin{proposition}[Memory loci in increment coordinates \eventproof]\label{prop:loci}
Define
\begin{equation}
 \Q_m=\{T\in\mathfrak T_N:\rank S_k(T)\le m\text{ for every }1\le k<N\}.
 \label{eq:memory-locus}
\end{equation}
Then
\begin{equation}
 s(W(d,T))=1+\max_k\rank S_k(T),\qquad
 \A_m\cong\R^N\times\Q_m.
 \label{eq:rank-increments}
\end{equation}
The locus $\Q_m$ is the lower quasiseparable-order $\le m$ locus \cite{EidelmanGohberg05}: it consists of those $T$ for which every $(m+1)$-minor of $S_k(T)$ equals zero for every $k=1,\ldots,N-1$. In particular, $\Q_0=\{0\}$, and the one-register maps are precisely
\begin{equation}
 y_{i+1}=y_i+d_i\nabla f(y_i),\qquad i=0,\ldots,N-1.
 \label{eq:GD}
\end{equation}
\end{proposition}
\begin{proof}
The coordinate formulas follow by telescoping. Retain the first row of $H_k$ and replace every subsequent row by its difference from the preceding original row. This invertible row operation gives a matrix with first column $(1,0,\ldots,0)^\top$ and lower-right block $S_k(T)$. Hence $\rank H_k=1+\rank S_k$. Theorem~\ref{thm:realization} and the elementary minor criterion prove~\eqref{eq:rank-increments}.

Every entry $T_{ij}$ belongs to $S_{j+1}(T)$, so $\Q_0=\{0\}$. Subtracting consecutive history equations gives
\[
 y_{i+1}=y_i+d_ig_i+\sum_{j<i}T_{ij}g_j.
\]
At $T=0$ this is~\eqref{eq:GD}, implemented by storing the current iterate. The converse follows from~\eqref{eq:rank-increments}.
\end{proof}

Let $\Phi_m(T)$ denote the conjunction that all the specified minors vanish. A block too small to have an $(m+1)$-minor contributes no equation. This convention also covers $N=1$.

\section{The three-step procedure and its correctness}
\label{sec:procedure}
For a real performance threshold $\gamma$, define the sets
\begin{equation}
 \begin{aligned}
 X_\gamma&=\{W\in\W:\risk(W)<\gamma\},\\
 P_\gamma&=\{T\in\mathfrak T_N:\exists d\in\R^N,\ W(d,T)\in X_\gamma\}.
 \end{aligned}
 \label{eq:performance}
\end{equation}
Thus $X_\gamma$ is a set of coefficient maps, and $P_\gamma$ is a set of increment matrices: it contains exactly those $T$ for which some choice of $d$ gives risk strictly below $\gamma$. Risk need not be constant on a fixed-$T$ orbit; the existential choice of $d$ is part of the performance problem.

\paragraph{Step 1: project the performance condition.}
Starting directly from Assumption~\ref{ass:finite}, write $P_\gamma$ explicitly as
\begin{equation}
 P_\gamma=\left\{T\in\mathfrak T_N\;\middle|\;
 \begin{array}{l}
 \exists d\in\R^N\ \exists b\in\R\text{ such that }b<\gamma\text{ and}\\[2pt]
 \forall z\in\R^q,\quad
 \Psi_N(W(d,T),z)\Rightarrow e_N(W(d,T),z)\le b
 \end{array}
 \right\}.
 \label{eq:step1}
\end{equation}
Eliminate $d,b,z$ over the real field to obtain a quantifier-free semialgebraic description of membership $T\in P_\gamma$, with $(T,\gamma)$ as the free variables.

\paragraph{Step 2: impose the memory equations.}
Construct the explicit polynomial equations $\Phi_m(T)$ from Proposition~\ref{prop:loci}, for $m=0,\ldots,\lfloor N/2\rfloor$.

\paragraph{Step 3: decide the intersections.}
For each such $m$, define the set of feasible performance thresholds
\begin{equation}
 \begin{aligned}
 \mathsf F_m&=\{\gamma\in\R:P_\gamma\cap\Q_m\ne\varnothing\}\\
 &=\{\gamma\in\R:\exists T\in\mathfrak T_N,\ T\in P_\gamma\ \land\ \Phi_m(T)\}.
 \end{aligned}
 \label{eq:step3}
\end{equation}
The elements of $\mathsf F_m$ are real thresholds, whereas the elements of $P_\gamma$ are matrices. Eliminate $T$ in~\eqref{eq:step3} to obtain a semialgebraic description of $\mathsf F_m\subseteq\R$. Isolate its finite boundary, if it has one, and compare the resulting endpoints between budgets.

\begin{theorem}[Correctness, termination, and exact output \eventproof]\label{thm:procedure}
Under Assumption~\ref{ass:finite}, the three steps terminate and satisfy
\begin{align}
 T\in P_\gamma&\Longleftrightarrow
 \exists d\in\R^N,\ \risk(W(d,T))<\gamma,\label{eq:P-correct}\\
 \gamma\in \mathsf F_m&\Longleftrightarrow
 P_\gamma\cap\Q_m\ne\varnothing
 \Longleftrightarrow X_\gamma\cap\A_m\ne\varnothing.\label{eq:F-correct}
\end{align}
Moreover,
\begin{equation}
 \mathsf F_m=
 \begin{cases}
 (V_m,\infty),&V_m<\infty,\\
 \varnothing,&V_m=\infty.
 \end{cases}
 \label{eq:ray}
\end{equation}
Consequently each finite $V_m$ is a computable real algebraic number. The procedure decides every strict memory gain and, when $V_{\mathrm{lin}}<\infty$, computes $m_\star$.
\end{theorem}
\begin{proof}
By the equality of loss sets in~\eqref{eq:loss-model}, the universally quantified condition in~\eqref{eq:step1} is exactly $\risk(W(d,T))\le b$. For a real $\gamma$, the inequality $\risk(W)<\gamma$ holds if and only if there is a real $b<\gamma$ with $\risk(W)\le b$. This proves the set equality~\eqref{eq:step1} and~\eqref{eq:P-correct}. The uniform bound $b<\gamma$ is needed even if the supremum loss is not attained.

Proposition~\ref{prop:loci} says that $\Phi_m(T)$ is precisely the memory condition. Thus a pair $(d,T)$ meeting both conditions exists exactly when $P_\gamma\cap\Q_m\ne\varnothing$, proving~\eqref{eq:F-correct}.

If $\gamma\in \mathsf F_m$, an eligible method has risk below $\gamma$, so $V_m<\gamma$. Conversely, if $V_m<\gamma$, the defining property of the infimum gives an eligible method with risk below $\gamma$. If $V_m=\infty$, no finite threshold works. This proves~\eqref{eq:ray} without assuming minimax attainment.

All formulas have finitely many real variables and polynomial conditions with real algebraic coefficients. Effective real quantifier elimination, followed by univariate real-root isolation, therefore performs the specified steps~\cite{BPR}. A finite endpoint of a semialgebraic subset of $\R$ defined over $\K$ is algebraic: after elimination, its boundary is among the real zeros of the finitely many nonzero polynomials in the resulting formula. Root isolation and sign determination identify the required endpoint exactly. Emptiness detects $V_m=\infty$.

Finally, Theorem~\ref{thm:realization} gives $V_{\lfloor N/2\rfloor}=V_{\mathrm{lin}}$. There are finitely many budgets to compare. Comparing their algebraic values, with the stated convention for infinity, decides every gain and gives the least sufficient budget when the unrestricted value is finite.
\end{proof}

\begin{corollary}[An exact threshold criterion for a memory gain \eventproof]\label{cor:gain}
For every $m\ge0$,
\begin{equation}
 V_{m+1}<V_m
 \quad\Longleftrightarrow\quad
 \mathsf F_{m+1}\setminus \mathsf F_m\ne\varnothing.
 \label{eq:gain}
\end{equation}
The set difference consists of thresholds attainable with $m+1$ auxiliary registers but not with $m$. If a gain exists, a threshold $\gamma\in \mathsf F_{m+1}\setminus \mathsf F_m$ can be chosen rational.
\end{corollary}
\begin{proof}
A threshold strictly between distinct values belongs to the set difference by~\eqref{eq:ray}; if the larger value is infinite, choose any sufficiently large finite threshold. Rational density allows a rational choice. Conversely any $\gamma\in \mathsf F_{m+1}\setminus \mathsf F_m$ gives $V_{m+1}<\gamma\le V_m$.
\end{proof}

The procedure is exact for each supplied finite horizon. It requires no minimax value or optimal algorithm as input. Its correctness does not imply a practical running-time bound, and algebraicity does not imply an expression by radicals.

\section{Extension to projected subgradient methods}
\label{sec:projected}
This section replaces the unconstrained gradient setting by a fixed nonempty closed convex set $K\subseteq\R^n$ and a subgradient oracle. The dimension $n$ is fixed here. For effective computation, $K$ is supplied by a finite semialgebraic membership formula $\kappa(x)$ with known real algebraic coefficients. Closedness ensures that the Euclidean projection $\Pi_K:\R^n\to K$ exists everywhere and is unique. Compactness is not required.

An instance consists of a finite convex function $f:K\to\R$, an initial point $x_0\in K$, and a minimizer $x_\star\in K$ with $\|x_0-x_\star\|\le R$. Define
\[
 \partial_K f(x)=\{g\in\R^n:
 f(v)\ge f(x)+\langle g,v-x\rangle\text{ for every }v\in K\}.
\]
Assume $\partial_K f(x)\ne\varnothing$ at every $x\in K$. Each oracle call returns one arbitrary $g\in\partial_K f(x)$. Risk includes all permitted choices, including different returns at repeated queries. Projection is an additional exact primitive: its internal workspace is not persistent algorithmic memory, and its calls are counted separately from subgradient queries.

\subsection{The projected recurrence}
Use the same coordinates $W\leftrightarrow(d,T)$ from~\eqref{eq:coordinates}--\eqref{eq:inverse}. Starting at $x_0$, perform, for $k=0,\ldots,N-1$,
\begin{align}
 g_k&\in\partial_K f(x_k),\nonumber\\
 u_{k+1}&=x_k+d_kg_k+\sum_{j<k}T_{kj}g_j,\label{eq:projected-increments}\\
 x_{k+1}&=\Pi_K(u_{k+1}).\nonumber
\end{align}
Return $x_N$. This uses $N$ subgradient queries and $N$ projections, including the final projection. All query points and the output belong to $K$. When $T=0$ and $d_k=-\alpha_k$, it is projected subgradient descent,
\[
 x_{k+1}=\Pi_K(x_k-\alpha_kg_k).
\]
When $K=\R^n$, it reduces to the original coefficient map $W$. The coefficients retain the same predetermined, instance-independent meaning as before; they may additionally depend on the known set $K$.

The projection is applied to each \emph{increment}, and its result is used in the next increment. In general this differs from projecting the original history expression $x_0+\sum_{j<i}a_{ij}g_j$ at each query. To record the difference exactly, define the projection corrections
\[
 \delta_i=x_i-u_i\qquad(1\le i\le N).
\]
Induction in~\eqref{eq:projected-increments} gives
\begin{equation}
 x_i=x_0+\sum_{j<i}a_{ij}g_j+\sum_{\ell=1}^{i}\delta_\ell.
 \label{eq:projected-history}
\end{equation}
Thus every past correction has coefficient one in every later iterate.

\subsection{Persistent registers and the exact memory formula}
\begin{definition}[Projected register realization]\label{def:projected-memory}
An $s$-register realization of~\eqref{eq:projected-increments} uses $Z_k\in(\R^n)^s$, initialized by $Z_0=e_1\otimes x_0$, and predetermined constants
\[
 F_k\in\R^{s\times s},\quad b_k,p_k\in\R^s\quad(0\le k<N),
 \qquad c_k\in\R^s\quad(0\le k\le N),
\]
with $c_0^\top e_1=1$ and $c_{k+1}^\top p_k=1$. It executes
\begin{align}
 x_k&=(c_k^\top\otimes I_n)Z_k,\nonumber\\
 \widehat Z_{k+1}&=(F_k\otimes I_n)Z_k+(b_k\otimes I_n)g_k,\nonumber\\
 \widehat u_{k+1}&=(c_{k+1}^\top\otimes I_n)\widehat Z_{k+1},\label{eq:projected-state}\\
 \widehat\delta_{k+1}&=\Pi_K(\widehat u_{k+1})-\widehat u_{k+1},\nonumber\\
 Z_{k+1}&=\widehat Z_{k+1}+(p_k\otimes I_n)\widehat\delta_{k+1}.\nonumber
\end{align}
The normalization ensures that the next observation is $\Pi_K(\widehat u_{k+1})$. Only $Z_{k+1}$ persists after the update. There is no additional stored transcript or free persistent copy of $x_0$.

The realization must reproduce every $x_i$ in~\eqref{eq:projected-increments} for every $x_0\in K$ and every sequence of formal vectors $g_0,\ldots,g_{N-1}\in\R^n$. These vectors need not be subgradients of a common function. Projection is still evaluated as the actual map $\Pi_K$. Define $s_K^\Pi(W)$ as the least such $s\ge1$.
\end{definition}

The formal-input requirement is the projected counterpart of Definition~\ref{def:memory}. It measures realization of the prescribed map, independently of which function class supplies the inputs. Equality on actual oracle executions follows immediately.

\begin{samepage}
\begin{theorem}[Cut rank with projection after every increment \eventproof]\label{thm:projected-rank}
If $K$ contains at least two points, then
\begin{equation}
 s_K^\Pi(W)=\max_{0\le k\le N}\rank H_k(W)
 =1+\max_{1\le k<N}\rank S_k(T).
 \label{eq:projected-rank}
\end{equation}
For $N=1$, the maximum over increment blocks is defined as zero. Consequently the maps admitting at most $m+1$ registers in Definition~\ref{def:projected-memory} are exactly
\[
 \A_m^\Pi\cong\R^N\times\Q_m.
\]
For a singleton $K$, every projected trajectory is constant and $s_K^\Pi(W)=1$ for every $W$.
\end{theorem}
\end{samepage}
\begin{proof}
To prove $s_K^\Pi(W)\le\max_k\rank H_k(W)$, put $E_k=\im H_k(W)$ as in Theorem~\ref{thm:realization}. Define the canonical state as a list of $n$-vectors with entries
\[
 (v_k)_i=x_0+\sum_{j<k}a_{ij}g_j+
                  \sum_{\ell=1}^k\delta_\ell,
 \qquad i=k,\ldots,N.
\]
The all-ones vector belongs to $E_k$, so $v_k\in E_k\otimes\R^n$. Its first $n$-vector is $x_k$. With $D_k$ and $h_k$ as in the earlier proof, form
\[
 w_{k+1}=(D_k\otimes I_n)v_k+h_k\otimes g_k
 \in E_{k+1}\otimes\R^n.
\]
Its first $n$-vector is $u_{k+1}$. Apply $\Pi_K$ to that vector, form $\delta_{k+1}$, and set
\[
 v_{k+1}=w_{k+1}+\mathbf 1\otimes\delta_{k+1}
 \in E_{k+1}\otimes\R^n.
\]
This gives exactly the next canonical state by~\eqref{eq:projected-history}. Choosing coordinates in each $E_k$ and padding to $s=\max_k\dim E_k$ gives~\eqref{eq:projected-state}. Explicitly, use the injections $J_k$ and left inverses $L_k$ from the earlier proof, take $p_k=J_{k+1}\mathbf 1$, and take $c_{k+1}^\top$ to be the first coordinate of $L_{k+1}$. Then $c_{k+1}^\top p_k=1$. Tensoring the linear maps with $I_n$ completes the construction.

To prove $s_K^\Pi(W)\ge\max_k\rank H_k(W)$, let an $s$-register realization be given. Set $L_K=\operatorname{span}(K-K)$ and choose a unit vector $v\in L_K$ and a point $\bar x\in\operatorname{relint}K$. Such a vector exists because $K$ is not a singleton. Choose formal inputs
\[
 x_0=\bar x+\alpha v,\qquad g_j=\beta_jv.
\]
For $(\alpha,\beta_0,\ldots,\beta_{N-1})$ in a sufficiently small open neighborhood of zero, every unprojected history point remains in a relative neighborhood of $\bar x$ contained in $\operatorname{relint}K$. Thus the prescribed trajectory has inactive projections and
\[
 x_i=\bar x+\left(\alpha+\sum_{j<i}a_{ij}\beta_j\right)v.
\]
The realizing program has the same projected observations. Whenever $\Pi_K(u)=x\in\operatorname{relint}K$, projection optimality implies $u-x\perp L_K$: test the projection inequality against $x\pm\varepsilon w\in K$ for $w\in L_K$ and sufficiently small $\varepsilon>0$. In particular,
\[
 \langle v,\widehat\delta_{k+1}\rangle=0.
\]
Taking the inner product with $v$ in every register therefore removes the correction term from~\eqref{eq:projected-state}. The resulting scalar state follows the linear recurrence with coefficients $F_k,b_k,c_k$ and reproduces the scalar history map $W$ on an open set of scalar initial values and inputs. Two linear maps agreeing on an open set agree everywhere. This is an $s$-register linear realization of $W$, so Theorem~\ref{thm:realization} gives the lower bound. Proposition~\ref{prop:loci} gives the increment formula.

For $K=\{a\}$, the one-register program retaining $x_0=a$ reproduces every projected trajectory.
\end{proof}

The theorem concerns exactly Definition~\ref{def:projected-memory} and the family~\eqref{eq:projected-increments}. It does not characterize memory for arbitrary nonlinear state encodings or unrestricted additional projection operations.

\subsection{An exact finite semialgebraic loss model}
For the supplied membership formula $\kappa$, define the projection relation
\begin{equation}
 \operatorname{Proj}_K(u,x) \Longleftrightarrow\
 \kappa(x)\ \land\
 \forall v\in\R^n\,
 [\,\kappa(v)\Rightarrow\langle u-x,v-x\rangle\le0\,].
 \label{eq:projection-graph}
\end{equation}
\begin{lemma}[Exact semialgebraic projection graph \eventproof]\label{lem:projection-graph}
Relation~\eqref{eq:projection-graph} holds if and only if $x=\Pi_K(u)$. Its graph has an effectively computable quantifier-free semialgebraic description.
\end{lemma}
\begin{proof}
If $x$ minimizes squared distance to $u$, comparison with $x+\varepsilon(v-x)\in K$ and passage to $\varepsilon\downarrow0$ gives the displayed inequality. Conversely, that inequality implies, for every $v\in K$,
\[
 \|u-v\|^2
 =\|u-x\|^2+\|v-x\|^2-2\langle u-x,v-x\rangle
 \ge\|u-x\|^2+\|v-x\|^2.
\]
Hence $x$ is the unique projection. The formula has finitely many real variables and polynomial conditions over $\K$, so effective real quantifier elimination removes $v$~\cite{BPR}.
\end{proof}

Let $J=\{\star,0,\ldots,N\}$. The loss variables $z$ consist of points $x_i$ for $i\in J$, vectors $g_0,\ldots,g_N$, scalars $\varphi_0,\ldots,\varphi_N$, and preprojection points $u_1,\ldots,u_N$. Set $g_\star=0$ and $\varphi_\star=0$. Define $\Psi_{N,K}^\Pi(W,z)$ by the following finite conditions:
\begin{align}
 &x_i\in K\quad(i\in J),\qquad
   \|x_0-x_\star\|^2\le R^2,\label{eq:projected-loss-model}\\
 &\varphi_i\ge\varphi_j+\langle g_j,x_i-x_j\rangle
                      \quad(i,j\in J),\nonumber\\
 &u_{k+1}=x_k+d_kg_k+\sum_{j<k}T_{kj}g_j,
                         \quad 0\le k<N,\nonumber\\
 &\operatorname{Proj}_K(u_{k+1},x_{k+1})
                         \quad 0\le k<N.\nonumber
\end{align}
Use the quantifier-free version of~\eqref{eq:projection-graph} in this definition. The represented loss is $e_{N,K}^\Pi(W,z)=\varphi_N$. The vector $g_N$ is an auxiliary terminal subgradient; it does not require another oracle query.

\begin{proposition}[Exactness of the projected loss description \eventproof]\label{prop:projected-loss}
For every fixed $W$, the feasible values of $\varphi_N$ in~\eqref{eq:projected-loss-model} are exactly the attainable losses of~\eqref{eq:projected-increments} over the stated instance class and permitted oracle choices.
\end{proposition}
\begin{proof}
Every actual trajectory supplies feasible data by the subgradient inequalities and Lemma~\ref{lem:projection-graph}, after subtracting $f(x_\star)$ from all function values. The terminal subgradient exists by assumption. At a constrained minimizer, $0\in\partial_K f(x_\star)$ follows directly from the definition.

Conversely, from feasible data construct
\begin{equation}
 f(x)=\max_{j\in J}
       \{\varphi_j+\langle g_j,x-x_j\rangle\},\qquad x\in K.
 \label{eq:projected-interpolant}
\end{equation}
The pairwise inequalities imply that $f(x_i)=\varphi_i$ and that $g_i\in\partial_K f(x_i)$ for every $i\in J$. This finite maximum is convex and has a subgradient everywhere: an active affine term supplies one. The $\star$-term is zero, so $f\ge0=f(x_\star)$. The remaining conditions and uniqueness of projection reproduce the specified trajectory by induction. Thus every feasible loss is attained. This is the standard convex interpolation construction underlying performance estimation, including projected methods~\cite{THGcomp}.
\end{proof}

Unbounded subgradients need not give finite minimax values. A precise bounded variant fixes $M>0$, requires at least one element of $\partial_K f(x)$ of norm at most $M$ at every $x\in K$, and restricts each oracle return to such elements. This does not require all relative subgradients at boundary points to be bounded. Add
\begin{equation}
 \|g_i\|^2\le M^2\qquad(0\le i\le N)
 \label{eq:bounded-projected-subgradients}
\end{equation}
to~\eqref{eq:projected-loss-model}. Proposition~\ref{prop:projected-loss} remains exact: every slope in~\eqref{eq:projected-interpolant} has norm at most $M$, so every point has a permitted active subgradient. The zero-step method has risk at most $MR$, since the subgradient inequality at $x_0$ gives
\[
 f(x_0)-f(x_\star)
 \le\langle g_0,x_0-x_\star\rangle\le MR.
\]
Consequently every memory-constrained minimax value is finite in this bounded variant.

\subsection{The three steps with projections}
Write $\risk_K^\Pi(W)$ for the supremum of $\varphi_N$ in the chosen version of~\eqref{eq:projected-loss-model}. For a non-singleton $K$, define
\begin{equation}
 P_\gamma^\Pi=\left\{T\in\mathfrak T_N\;\middle|\;
 \begin{array}{l}
 \exists d\in\R^N\ \exists b\in\R\text{ such that }b<\gamma\text{ and}\\[2pt]
 \forall z,\quad
 \Psi_{N,K}^\Pi(W(d,T),z)\Rightarrow\varphi_N\le b
 \end{array}
 \right\},
 \label{eq:projected-performance}
\end{equation}
and
\begin{equation}
 \mathsf F_m^\Pi
 =\{\gamma\in\R:P_\gamma^\Pi\cap\Q_m\ne\varnothing\}.
 \label{eq:projected-thresholds}
\end{equation}
All variables $z$ range over their indicated real coordinate spaces. The three steps are to eliminate $d,b,z$ in~\eqref{eq:projected-performance}, impose the same minor equations defining $\Q_m$, and eliminate $T$ in~\eqref{eq:projected-thresholds}.

\begin{samepage}
\begin{theorem}[Exact projected memory decision procedure \eventproof]\label{thm:projected-procedure}
Fix $N,n$, a non-singleton closed convex semialgebraic $K$ described over $\K$, and $R\in\K$ with $R>0$. For the bounded variant also fix $M\in\K$ with $M>0$. Define
\[
 V_m^\Pi
 =\inf_{\{W:s_K^\Pi(W)\le m+1\}}\risk_K^\Pi(W).
\]
The projected three-step procedure terminates, and
\begin{equation}
 \mathsf F_m^\Pi=
 \begin{cases}
 (V_m^\Pi,\infty),&V_m^\Pi<\infty,\\
 \varnothing,&V_m^\Pi=\infty.
 \end{cases}
 \label{eq:projected-ray}
\end{equation}
It computes every finite $V_m^\Pi$ as a real algebraic number and decides every strict comparison $V_{m+1}^\Pi<V_m^\Pi$. In the bounded variant all values are finite, and the procedure computes the least auxiliary budget achieving the unrestricted value within family~\eqref{eq:projected-increments}.
\end{theorem}
\end{samepage}
\begin{proof}
Proposition~\ref{prop:projected-loss} identifies the supremum in the finite model with actual risk. Thus~\eqref{eq:projected-performance} describes precisely the increment matrices admitting some $d$ with risk below $\gamma$. Theorem~\ref{thm:projected-rank} identifies $\Q_m$ with the required memory bound. Their intersection is nonempty exactly when an eligible method has risk below $\gamma$. The infimum argument in Theorem~\ref{thm:procedure} proves~\eqref{eq:projected-ray}. All conditions are finite semialgebraic formulas over $\K$, so the same quantifier-elimination and root-isolation argument proves effective computation and exact comparisons. The rank bound again makes $m=\lfloor N/2\rfloor$ unrestricted. Finiteness in the bounded variant follows from the zero-step bound.
\end{proof}

This result concerns a specified domain in a fixed finite dimension. A supremum over domains in unbounded dimensions would require an additional uniform finite reduction. It is not asserted here.

\FloatBarrier

The next two sections return to the principal unconstrained setting of Section~\ref{sec:problem}, including all finite dimensions, and compute its memory thresholds.

\section{Computation for convex quadratics}
\label{sec:quadratics}
Fix $L,R>0$ and consider, in every finite dimension,
\begin{equation}
 f(x)=\tfrac12(x-x_\star)^\top Q(x-x_\star)+c,\qquad
 Q=Q^\top,\quad 0\preceq Q\preceq LI,\quad \|x_0-x_\star\|\le R.
 \label{eq:quadratic-class}
\end{equation}
For the effective claims take $L,R\in\K$; the explicit formulas below hold for arbitrary positive real $L,R$.

\subsection{Step 1: construct the exact performance formula}
For a fixed $W=W(d,T)$ define scalar polynomials recursively, for $0\le i<N$, by
\begin{equation}
 p_0(\lambda)=1,\qquad
 p_{i+1}(\lambda)=(1+d_i\lambda)p_i(\lambda)
                 +\lambda\sum_{j<i}T_{ij}p_j(\lambda).
 \label{eq:polynomials}
\end{equation}
Their coefficients are polynomial functions of $(d,T)$, with
$\deg p_i\le i$ and $p_i(0)=1$.

For a polynomial $p_i(\lambda)=\sum_{\ell=0}^i c_{i\ell}\lambda^\ell$, its evaluation at the matrix $Q\in\R^{n\times n}$ means
\begin{equation}
 p_i(Q):=\sum_{\ell=0}^i c_{i\ell}Q^\ell,
 \qquad Q^0=I_n.
 \label{eq:matrix-polynomial}
\end{equation}
This is an $n\times n$ matrix, obtained by replacing each scalar power $\lambda^\ell$ by the matrix power $Q^\ell$. In particular, $p_0(Q)=I_n$.

Write $r_i=y_i-x_\star$ for the error vector at iterate $i$. The polynomial $p_i$ describes how the initial error $r_0=x_0-x_\star$ propagates through the method.

\begin{proposition}[Quadratic error and loss formulas \eventproof]\label{prop:quadratic-risk}
For every coefficient map, the error-propagation identity is
\begin{equation}
 r_i=p_i(Q)r_0,
 \quad\text{that is,}\quad
 y_i-x_\star=p_i(Q)(x_0-x_\star)
 \qquad(0\le i\le N).
 \label{eq:quad-error}
\end{equation}
The resulting risk is
\begin{equation}
 \risk(W)=\frac{R^2}{2}\max_{0\le\lambda\le L}\lambda p_N(\lambda)^2.
 \label{eq:quad-risk}
\end{equation}
An exact finite loss model is
\[
 0\le\lambda\le L,\qquad 0\le\rho\le R^2,\qquad
 e_N(W,\lambda,\rho)=\tfrac{\rho}{2}\lambda p_N(\lambda)^2.
\]
In particular the projected performance set is
\begin{equation}
 P_\gamma=\left\{T\in\mathfrak T_N\;\middle|\;
 \begin{array}{l}
 \exists d\in\R^N\ \exists b\in\R\text{ such that }b<\gamma\text{ and}\\[2pt]
 \forall\lambda\in[0,L],\quad
       \tfrac{R^2}{2}\lambda p_N(\lambda;d,T)^2\le b
 \end{array}
 \right\}.
 \label{eq:quad-projection}
\end{equation}
\end{proposition}
\begin{proof}
For the quadratic~\eqref{eq:quadratic-class}, $g_j=Q(y_j-x_\star)=Qr_j$. Subtracting $x_\star$ from the increment recurrence gives
\[
 r_{i+1}=(I_n+d_iQ)r_i+Q\sum_{j<i}T_{ij}r_j.
\]
Identity~\eqref{eq:quad-error} holds at $i=0$ because $p_0(Q)=I_n$. If it holds through index $i$, then
\[
 r_{i+1}
 =\left[(I_n+d_iQ)p_i(Q)+Q\sum_{j<i}T_{ij}p_j(Q)\right]r_0
 =p_{i+1}(Q)r_0,
\]
where the last equality is the polynomial recurrence~\eqref{eq:polynomials} evaluated at $Q$. This proves the error-propagation identity by induction.

To obtain the risk, diagonalize $Q=U\operatorname{diag}(\lambda_1,\ldots,\lambda_n)U^\top$ with $U$ orthogonal and $0\le\lambda_a\le L$. The definition~\eqref{eq:matrix-polynomial} gives
\[
 p_N(Q)=U\operatorname{diag}(p_N(\lambda_1),\ldots,p_N(\lambda_n))U^\top.
\]
Writing $\xi=U^\top r_0$, we obtain
\[
 f(y_N)-f(x_\star)
 =\frac12\sum_{a=1}^n\lambda_a p_N(\lambda_a)^2\xi_a^2
 \le\frac{R^2}{2}\max_{0\le\lambda\le L}\lambda p_N(\lambda)^2.
\]
A scalar quadratic with a maximizing eigenvalue and initial error $R$ attains this bound, proving~\eqref{eq:quad-risk}.

Let $M_W=\max_{[0,L]}\lambda p_N(\lambda)^2$. Every multidimensional loss lies in $[0,R^2M_W/2]$. Every value in that interval is attained by a scalar quadratic at a maximizing eigenvalue, with initial squared error $\rho\in[0,R^2]$. Hence the finite model describes exactly the attainable losses, not merely their supremum. In the uniform upper-bound condition from Step 1, $\rho$ can be eliminated by taking $\rho=R^2$, since $\lambda p_N(\lambda)^2\ge0$. This gives precisely~\eqref{eq:quad-projection}.
\end{proof}

Formula~\eqref{eq:quad-projection} is the quadratic specialization of~\eqref{eq:step1}; it is still quantified in $d,b,\lambda$. The general algorithm would eliminate these variables to describe all of $P_\gamma$. Here we compute the final threshold sets $\mathsf F_m$ directly by the analytic argument below, combining the eliminations in Steps 1 and 3. We do not compute a quantifier-free description of $P_\gamma$ for arbitrary $T$.

\subsection{Step 2: impose the memory equations}
For each budget $m$, use exactly the equations $\Phi_m(T)$ from Section~\ref{sec:memory}: every $(m+1)$-minor of every block $S_k(T)$ must equal zero, for $k=1,\ldots,N-1$. The function class does not change these memory equations.

For $m=0$, the minors are individual entries. Every entry $T_{ij}$ occurs in $S_{j+1}(T)$, so these equations are precisely
\[
 T_{ij}=0\quad(0\le j<i<N),\qquad\text{equivalently}\qquad\Q_0=\{0\}.
\]
Setting $T=0$ in~\eqref{eq:polynomials} gives the error polynomial
\begin{equation}
 p_N(\lambda;d,0)=\prod_{i=0}^{N-1}(1+d_i\lambda).
 \label{eq:product}
\end{equation}
The vector $d$ remains free: Step 2 constrains memory, while Step 3 must decide which choices of $d$ satisfy the performance requirement. Since $\Q_0\subseteq\Q_m$ for every $m\ge0$, a feasible method at $T=0$ will also be feasible for every larger memory budget.

\subsection{Step 3: compute the threshold sets}
The task is exactly the intersection test~\eqref{eq:step3}, namely to compute
\[
 \mathsf F_m=\{\gamma\in\R:P_\gamma\cap\Q_m\ne\varnothing\}.
\]
We prove nonemptiness above an explicit threshold by constructing a point $T=0$ in the intersection and coefficients $d$ witnessing membership in $P_\gamma$. We prove emptiness at or below that threshold by a lower bound valid for every $d,T$, hence for every memory budget. Together these give the quantifier-free description of $\mathsf F_m$ required by Step 3.

Let $M=2N+1$, and let $C_M$ be the Chebyshev polynomial specified by
$C_M(\cos\theta)=\cos(M\theta)$. The odd-Chebyshev construction appears in Vandenberghe \cite{Vandenberghe22}, Exercise~17; we use it here for the predetermined-step minimax problem. Define the polynomial $p_\star$ by
\begin{equation}
 u\,p_\star(Lu^2)=\frac{(-1)^N}{M}C_M(u).
 \label{eq:Cheb}
\end{equation}
Since $C_M$ is odd and $C_M'(0)=(-1)^NM$, this defines a degree-$N$ polynomial with $p_\star(0)=1$.

\begin{theorem}[Complete quadratic memory calculation \eventproof]\label{thm:quadratic}
For every $N\ge1$, every $m\ge0$, and every real $\gamma$,
\begin{equation}
 P_\gamma\cap\Q_m\ne\varnothing
 \quad\Longleftrightarrow\quad
 \gamma>\frac{LR^2}{2(2N+1)^2}.
 \label{eq:quad-intersection}
\end{equation}
Thus the output of Step 3 is the explicitly computed set
\begin{equation}
 \mathsf F_m=\left(\frac{LR^2}{2(2N+1)^2},\infty\right)
 \qquad(m\ge0).
 \label{eq:quad-thresholds}
\end{equation}
Taking its finite endpoint, as prescribed in Theorem~\ref{thm:procedure}, gives
\begin{equation}
 V_m=V_{\mathrm{lin}}=\frac{LR^2}{2(2N+1)^2}
 \quad(m\ge0),\qquad m_\star=0.
 \label{eq:quad-values}
\end{equation}
The value is attained by one-register scheduled gradient descent.
\end{theorem}
\begin{proof}
On $[0,1]$, $|C_M(u)|\le1$, so~\eqref{eq:Cheb} gives
$\max_{[0,L]}\lambda p_\star(\lambda)^2=L/M^2$, with equality at $\lambda=L$.
The roots of $p_\star$ are the distinct positive numbers
\[
 \lambda_j=L\cos^2\left(\frac{(2j-1)\pi}{2M}\right),
 \qquad 1\le j\le N.
\]
Its constant term is one, hence
$p_\star(\lambda)=\prod_{j=1}^N(1-\lambda/\lambda_j)$.
Set $d_{j-1}=-1/\lambda_j$ and $T=0$. Equations~\eqref{eq:GD} and~\eqref{eq:product} give an explicit one-register implementation with risk $LR^2/(2M^2)$.
More explicitly, if $\gamma>LR^2/(2M^2)$, choose
\[
 b=\frac12\left(\gamma+\frac{LR^2}{2M^2}\right).
\]
Then $b<\gamma$ and the constructed coefficients satisfy the uniform bound in~\eqref{eq:quad-projection}. Hence $0\in P_\gamma$, and Step 2 gives $0\in\Q_m$ for every $m$. This proves nonemptiness of the required intersections.

For the lower bound, suppose some polynomial $p$ of degree at most $N$ with $p(0)=1$ satisfied $\max_{0\le\lambda\le L}\lambda p(\lambda)^2<L/M^2$. Put
$u_i=\cos(i\pi/M)>0$, $i=0,\ldots,N$.
Then $u_i p_\star(Lu_i^2)=(-1)^{N+i}/M$, whereas
$|u_i p(Lu_i^2)|<1/M$. The difference $p-p_\star$ therefore has alternating strict signs at the $N+1$ distinct positive points $Lu_i^2$. It has at least $N$ distinct roots between them and also a root at zero. This contradicts its degree bound.

Every predetermined method $W(d,T)$ produces a polynomial $p_N$ of degree at most $N$ with $p_N(0)=1$. Proposition~\ref{prop:quadratic-risk} therefore gives $\risk(W(d,T))\ge LR^2/(2M^2)$ for every $d,T$. If $\gamma\le LR^2/(2M^2)$, no such pair can meet the strict performance requirement, so $P_\gamma=\varnothing$ and $P_\gamma\cap\Q_m=\varnothing$ for every $m$. This proves the other direction of~\eqref{eq:quad-intersection}. Formula~\eqref{eq:quad-thresholds} and all minimax conclusions follow.
\end{proof}

The root schedule uses the known terminal horizon and unrestricted exact steps. It asserts a terminal guarantee; no constraint on intermediate objective values is part of the problem.

\section{Computation for convex smooth functions}
\label{sec:smooth}
Let $\mathcal F_{0,L}$ consist of convex continuously differentiable functions on $\R^n$ satisfying
\begin{equation}
 \|\nabla f(x)-\nabla f(y)\|\le L\|x-y\|\qquad(x,y\in\R^n),
 \label{eq:smooth-class}
\end{equation}
with an attained minimum and $\|x_0-x_\star\|\le R$. We again include every finite dimension and take $L,R>0$, with $L,R\in\K$ for the effective computation claims. Convexity is an explicit assumption in this example.

\subsection{Step 1: construct the exact performance formula}
We use the following established interpolation theorem, in its smooth convex specialization~\cite[Corollary 1]{THG}: finite data $(x_i,g_i,f_i)$ admit an interpolant in $\mathcal F_{0,L}$ if and only if
\begin{equation}
 f_i-f_j-\langle g_j,x_i-x_j\rangle
 -\frac1{2L}\|g_i-g_j\|^2\ge0\qquad\text{for every }i,j.
 \label{eq:interpolation}
\end{equation}
A datum $(x_\star,0,f_\star)$ enforces a global minimizer. Here an interpolant means one function whose value and gradient at every listed point equal the prescribed data; pairwise inequalities are therefore an exact consistency test for a single function.

To encode these inequalities independently of the instance dimension, put $\nu=N+2$ and introduce the list of $\nu$ physical vectors
\begin{equation}
 B=[\,x_0-x_\star\ \ g_0\ \ \cdots\ \ g_N\,]\in\R^{n\times\nu},
 \qquad G=B^\top B\in\R^{\nu\times\nu}.
 \label{eq:Gram-definition}
\end{equation}
Thus $G_{ab}$ is the inner product of columns $a$ and $b$ of $B$, indexed from zero. The vector $g_N=\nabla f(y_N)$ is used only for terminal interpolation; it is not an additional algorithmic query.

Let $e_0,\ldots,e_{N+1}$ be the standard basis of $\R^\nu$, and let $J=\{\star,0,\ldots,N\}$, with $y_\star=x_\star$ and $g_\star=0$. For $0\le i\le N$, define the coefficient vectors
\begin{equation}
 v_i(W)=e_0+\sum_{j<i}a_{ij}e_{j+1},\qquad w_i=e_{i+1},
 \qquad v_\star=w_\star=0.
 \label{eq:Gram-vectors}
\end{equation}
They live in $\R^\nu$, whereas the iterates and gradients live in $\R^n$. Their meaning is exactly
\begin{equation}
 Bv_i(W)=y_i-x_\star,\qquad Bw_i=g_i.
 \label{eq:Gram-representation}
\end{equation}
For arbitrary coefficient vectors $u,v\in\R^\nu$, the Gram identity is
$\langle Bu,Bv\rangle=u^\top Gv$. In particular,
\begin{align}
 \langle g_j,y_i-y_j\rangle
 &=w_j^\top G(v_i(W)-v_j(W)),\nonumber\\
 \|g_i-g_j\|^2&=(w_i-w_j)^\top G(w_i-w_j).
 \label{eq:Gram-products}
\end{align}

The finite model treats $G\in\Sym_\nu(\R)$ and $\varphi=(\varphi_0,\ldots,\varphi_N)\in\R^{N+1}$ as variables, where $\Sym_\nu(\R)$ is the space of real symmetric $\nu\times\nu$ matrices. The intended values are $\varphi_i=f(y_i)-f(x_\star)$, and $\varphi_\star=0$. Substituting~\eqref{eq:Gram-products} into~\eqref{eq:interpolation} gives the feasibility condition $\Psi_N^{\mathrm{sm}}(W,G,\varphi)$:
\begin{align}
 G&\succeq0,\qquad G_{00}\le R^2,\label{eq:Gram-constraints}\\
 \varphi_i-\varphi_j
 &-w_j^\top G(v_i(W)-v_j(W))
 -\frac1{2L}(w_i-w_j)^\top G(w_i-w_j)\ge0
 \quad(i,j\in J).\nonumber
\end{align}
Positive semidefiniteness is a finite polynomial condition: all principal minors are nonnegative. Also $G_{00}=\|x_0-x_\star\|^2$, so the radius condition is exactly the displayed bound. The loss is $e_N=\varphi_N$. The vector $z$ of Assumption~\ref{ass:finite} consists of the independent entries of $G$ and the $N+1$ function values, giving
\[
 q=\frac{\nu(\nu+1)}2+N+1
\]
real analysis variables. Neither $G$ nor these coefficient vectors and function values are stored by the algorithm.

\begin{proposition}[Exactness across all dimensions \eventproof]\label{prop:smooth-model}
The formula $\Psi_N^{\mathrm{sm}}$ and loss $\varphi_N$ satisfy Assumption~\ref{ass:finite} when $L,R$ are algebraic. Every attainable loss has a realization in dimension at most $N+2$. Consequently Step 1 gives the set
\begin{equation}
 P_\gamma=\left\{T\in\mathfrak T_N\;\middle|\;
 \begin{array}{l}
 \exists d\in\R^N\ \exists b\in\R\text{ such that }b<\gamma\text{ and}\\[2pt]
 \forall(G,\varphi)\in\Sym_\nu(\R)\times\R^{N+1},\\[2pt]
 \qquad\Psi_N^{\mathrm{sm}}(W(d,T),G,\varphi)\Rightarrow\varphi_N\le b
 \end{array}
 \right\}.
 \label{eq:smooth-projection}
\end{equation}
\end{proposition}
\begin{proof}
For an actual instance, form $B$ and $G=B^\top B$ as in~\eqref{eq:Gram-definition}, and set $\varphi_i=f(y_i)-f(x_\star)$. Identities~\eqref{eq:Gram-representation}--\eqref{eq:Gram-products} turn the interpolation inequalities into~\eqref{eq:Gram-constraints}. The terminal loss is $\varphi_N$. This represents every actual loss, whatever the original dimension $n$.

Conversely, positive semidefiniteness permits a factorization $G=B^\top B$ with $B\in\R^{\nu\times\nu}$. In dimension $\nu$, set $x_\star=0$, $x_0=Be_0$, $y_i=Bv_i(W)$, and $g_i=Bw_i$. Assign value $\varphi_i$ at $y_i$, and add the minimizer datum $(0,0,0)$. The same identities show that these data satisfy~\eqref{eq:interpolation}, so the cited theorem supplies a convex $L$-smooth interpolating function. Its prescribed gradients generate exactly the history recurrence~\eqref{eq:history}, and $G_{00}\le R^2$ enforces the initial radius. Thus every feasible loss is attained in dimension at most $\nu=N+2$. Substituting this exact loss model into~\eqref{eq:step1} gives~\eqref{eq:smooth-projection}.
\end{proof}

Formula~\eqref{eq:smooth-projection} is the quantified input to Step 1, with $d,b,G,\varphi$ to eliminate. As in Section~\ref{sec:quadratics}, the calculations below do not produce a quantifier-free description of all of $P_\gamma$. Instead, the cited optimality results and explicit memory certificates determine the required threshold intersections directly. For the one-register value at $N\ge2$, a final quantifier elimination remains to be performed if its exact numerical value is wanted.

\subsection{Step 2: impose the memory equations}
Use the same sets $\Q_m$ from~\eqref{eq:memory-locus}. For every $k=1,\ldots,N-1$, all $(m+1)$-minors of $S_k(T)$ must vanish. In particular,
\[
 \Q_0=\{0\},\qquad
 \Q_1=\{T\in\mathfrak T_N:\text{every }2\times2\text{ minor of every }S_k(T)\text{ is zero}\}.
\]
Thus $T=0$ is the one-register restriction, and $T\in\Q_1$ is the at-most-two-register restriction. The performance constraints from Step 1 are unchanged when a memory budget is imposed. Step 3 tests whether they can be met on these loci; the two-register construction below supplies a concrete point in $\Q_1$.

\subsection{Step 3: compute the thresholds for two or more registers}
The set to compute is again
\[
 \mathsf F_m=\{\gamma\in\R:P_\gamma\cap\Q_m\ne\varnothing\}.
\]
Define
\begin{equation}
 \theta_0=1,\qquad
 \theta_k=\frac{1+\sqrt{1+4\theta_{k-1}^2}}2\ (1\le k<N),\qquad
 \theta_N=\frac{1+\sqrt{1+8\theta_{N-1}^2}}2,
 \label{eq:theta}
\end{equation}
and put
\begin{equation}
 c_N=\frac{LR^2}{2\theta_N^2}.
 \label{eq:smooth-value}
\end{equation}
We invoke two established optimization results: the optimized gradient method attains the upper bound $c_N$~\cite[Theorem 2 and Section 7]{KF}, and every deterministic first-order method has worst-case loss at least $c_N$ in dimension $n\ge N+1$~\cite[Theorem 3 and Corollary 4]{Drori}. The latter applies in particular to Definition~\ref{def:methods}.

The numbers $\theta_k$, $\beta_k=1/\theta_{k+1}$, and $c_N$ depend only on the known horizon and class parameters. To construct the method, store the query vector $y_k$ and an auxiliary vector $\eta_k\in\R^n$, initially $(y_0,\eta_0)=(x_0,0)$. For $k=0,\ldots,N-1$, execute
\begin{align}
 g_k&=\nabla f(y_k),\nonumber\\
 y_{k+1}&=y_k+\beta_k\eta_k
             -\frac{1+(2\theta_k-1)\beta_k}{L}g_k,\label{eq:OGM-state}\\
 \eta_{k+1}&=(1-\beta_k)\eta_k
             -\frac{(2\theta_k-1)(1-\beta_k)}{L}g_k.\nonumber
\end{align}
Return $y_N$. Unrolling this linear recurrence defines the coefficient map $W_{\mathrm O}$ and hence its coordinates $(d_{\mathrm O},T_{\mathrm O})$ through~\eqref{eq:coordinates}. Indeed, $\eta_k$ is a linear combination of past gradients, while $y_k$ is $x_0$ plus such a combination. In particular,
\[
 (d_{\mathrm O})_k
 =-\frac{1+(2\theta_k-1)/\theta_{k+1}}{L}.
\]

\begin{proposition}[The OGM memory calculation \eventproof]\label{prop:OGM-memory}
The recurrence~\eqref{eq:OGM-state} is OGM and has risk $c_N$. Its coefficient map $W_{\mathrm O}$ satisfies
\[
 s(W_{\mathrm O})=
 \begin{cases}1,&N=1,\\2,&N\ge2.\end{cases}
\]
In particular, its increment matrix $T_{\mathrm O}$ belongs to $\Q_1$, and belongs to $\Q_0$ only when $N=1$.
\end{proposition}
\begin{proof}
For the identity with the published algorithm, put $z_k=y_k+\eta_k$. This is a quantity used in the proof. Adding the state updates gives
$z_{k+1}=z_k-2\theta_kg_k/L$, and the first update becomes
$y_{k+1}=(1-1/\theta_{k+1})(y_k-g_k/L)+z_{k+1}/\theta_{k+1}$.
These are the OGM2 equations~\cite[Section 7.2]{KF}, with $z_0=y_0=x_0$. The cited upper and lower bounds give its exact risk.

To check the memory condition explicitly, take $Z_k=(y_k,\eta_k)$ in Definition~\ref{def:memory}, with
\[
 F_k=\begin{pmatrix}1&\beta_k\\0&1-\beta_k\end{pmatrix},\qquad
 b_k=-\frac1L
 \begin{pmatrix}1+(2\theta_k-1)\beta_k\\(2\theta_k-1)(1-\beta_k)\end{pmatrix},
 \qquad c_k=o=e_1.
\]
This is a two-register realization initialized by $e_1\otimes x_0$. Proposition~\ref{prop:loci} therefore gives $\rank S_k(T_{\mathrm O})\le1$ at every cut: all the $2\times2$ minors in Step 2 vanish, so $T_{\mathrm O}\in\Q_1$.

For $N\ge2$, expansion of the first two outputs gives
\begin{equation}
 a_{10}=-\frac{1+1/\theta_1}{L},\qquad
 (T_{\mathrm O})_{10}=a_{20}-a_{10}
 =-\frac{1-1/\theta_1}{L\theta_2}\ne0.
 \label{eq:OGM-increment}
\end{equation}
This nonzero entry is a $1\times1$ minor in $S_1(T_{\mathrm O})$, so the one-register equations from Step 2 fail when $N\ge2$. For $N=1$, $\theta_1=2$ and the only output is $y_1=x_0-3g_0/(2L)$, which has a one-register implementation.
\end{proof}

\begin{theorem}[The smooth-convex plateau \eventproof; uses \literature]\label{thm:smooth-plateau}
For every $N\ge1$, $m\ge1$, and real $\gamma$,
\begin{equation}
 P_\gamma\cap\Q_m\ne\varnothing\quad\Longleftrightarrow\quad\gamma>c_N.
 \label{eq:smooth-intersections}
\end{equation}
Thus the quantifier-free output of Step 3 for these budgets is
\begin{equation}
 \mathsf F_m=(c_N,\infty)\qquad(m\ge1).
 \label{eq:smooth-thresholds}
\end{equation}
Its endpoint gives $V_m=V_{\mathrm{lin}}=c_N$ for all $m\ge1$. If $N=1$, the same statements hold for $m=0$.
\end{theorem}
\begin{proof}
If $\gamma\le c_N$, the cited lower bound gives $\risk(W(d,T))\ge c_N$ for every $d,T$. Therefore $P_\gamma=\varnothing$, and every intersection $P_\gamma\cap\Q_m$ is empty.

If $\gamma>c_N$, take $(d,T)=(d_{\mathrm O},T_{\mathrm O})$ and $b=(\gamma+c_N)/2$. Every feasible interpolation loss for this map is at most its risk $c_N<b<\gamma$. These are witnesses to membership $T_{\mathrm O}\in P_\gamma$ in~\eqref{eq:smooth-projection}. The memory calculation gives $T_{\mathrm O}\in\Q_1\subseteq\Q_m$, proving nonemptiness for every $m\ge1$. For $N=1$, it also gives $T_{\mathrm O}\in\Q_0$. These two cases prove the intersection criterion and the stated threshold sets. Their endpoints yield the minimax values by Theorem~\ref{thm:procedure}.
\end{proof}

\subsection{Step 3 continued: resolve the one-register budget}
For $N\ge2$, the remaining intersection is $P_\gamma\cap\Q_0=P_\gamma\cap\{0\}$. The previous construction does not settle it, since $T_{\mathrm O}\ne0$. Knowing that OGM requires two registers does not rule out a different one-register optimal method, or a sequence of one-register methods whose risks approach $c_N$. We use uniqueness to exclude the former and attainment to exclude the latter.

\begin{theorem}[External uniqueness theorem~{\cite[Theorem 1.1]{GJP}} \literature]\label{thm:uniqueness}
For $n\ge N+2$, OGM is the unique predetermined full-history coefficient map minimizing the worst-case terminal objective gap over convex $L$-smooth functions with the stated initial radius.
\end{theorem}
The cited result is from Grimmer, Jo, and Park, arXiv:2609.07671v1, Theorem~1.1 (September 2026). It fixes $L=1$ and uses the normalization $\tfrac12\|x_0-x_\star\|^2\le1$. For a general initial radius, the worst-case value scales by $R^2/2$ relative to the source normalization, while the minimizing coefficients transform as $a_{ij}=-W^{\mathrm{there}}_{ij}/L$. We use only the stated hypotheses. The strict separation below uses this theorem; Theorem~\ref{thm:smooth-plateau} does not.

\begin{lemma}[Attainment among one-register schedules \eventproof]\label{lem:compact}
For the smooth-convex class, the infimum $V_0$ is attained.
\end{lemma}
\begin{proof}
Write the one-register schedule as $d_i=-\alpha_i$, where $\alpha\in\R^N$, and abbreviate $\risk_0(\alpha)=\risk(W(-\alpha,0))$. Its error polynomial on the quadratic subclass, from~\eqref{eq:product}, is
$p_\alpha(\lambda)=\prod_{i=0}^{N-1}(1-\alpha_i\lambda)$.
These quadratics belong to the smooth-convex class, so~\eqref{eq:quad-risk} gives
\[
 \risk_0(\alpha)\ge\frac{R^2}{2}\lambda p_\alpha(\lambda)^2
 \qquad(0\le\lambda\le L).
\]
Consequently a risk bound $\risk_0(\alpha)\le C<\infty$ implies
\[
 |p_\alpha(\lambda)|\le\sqrt{\frac{4C}{LR^2}}
 \qquad(L/2\le\lambda\le L),
\]
since $\lambda\ge L/2$. To bound the steps themselves, evaluate $p_\alpha$ at $N+1$ fixed distinct points $\lambda_0,\ldots,\lambda_N\in[L/2,L]$. The coefficient vector is obtained from these values by inverting the fixed Vandermonde matrix with entries $A_{\ell j}=\lambda_\ell^j$. Thus all coefficients in $p_\alpha(\lambda)=1+\sum_{j=1}^N\zeta_j\lambda^j$ are bounded uniformly on the risk sublevel. The reversed monic polynomial is
\[
 \widehat p_\alpha(\xi)=\xi^N+\zeta_1\xi^{N-1}+\cdots+\zeta_N
 =\prod_{i=0}^{N-1}(\xi-\alpha_i).
\]
Its roots are precisely the steps $\alpha_i$. The elementary bound
$|\alpha_i|\le1+\max_{1\le j\le N}|\zeta_j|$ for roots of a monic polynomial therefore bounds every step uniformly.

For each fixed smooth instance, its terminal loss is continuous in $\alpha$ by the finite recurrence and continuity of the gradient and objective. Their supremum $\risk_0$ is lower semicontinuous. Hence each finite risk sublevel $\{\alpha\in\R^N:\risk_0(\alpha)\le C\}$ is closed as well as bounded, and is therefore compact.

The zero-step method has risk at most $LR^2/2$, by the smoothness inequality at a minimizer. A minimizing sequence therefore lies eventually in a nonempty finite sublevel. Compactness and lower semicontinuity produce an attaining schedule.
\end{proof}

\begin{samepage}
\begin{theorem}[Complete smooth-convex memory comparison \eventproof; uses \literature]\label{thm:smooth-memory}
Using Theorem~\ref{thm:uniqueness}, the least sufficient auxiliary budget is
\begin{equation}
 m_\star=
 \begin{cases}
 0,&N=1,\\
 1,&N\ge2.
 \end{cases}
 \label{eq:smooth-minmemory}
\end{equation}
For $N\ge2$ the only strict adjacent memory gain is
\begin{equation}
 V_0>V_1=V_2=\cdots=c_N.
 \label{eq:smooth-gap}
\end{equation}
\end{theorem}
\end{samepage}
\begin{proof}
The plateau and the case $N=1$ follow from Theorem~\ref{thm:smooth-plateau}. Suppose $N\ge2$ and $V_0=c_N$. By Lemma~\ref{lem:compact}, a one-register map $W_0$ attains that risk over all dimensions. In a fixed dimension $n\ge N+2$ its worst-case loss is at most $c_N$ and at least $c_N$ by the lower bound. It is therefore an optimal map in that dimension. Theorem~\ref{thm:uniqueness} forces $W_0=W_{\mathrm O}$. This contradicts~\eqref{eq:OGM-increment} and $\Q_0=\{0\}$. Thus $V_0>c_N$.
\end{proof}

To display exactly what remains of the general procedure, impose the Step 2 equation $T=0$ in~\eqref{eq:smooth-projection}. Step 3 then asks to eliminate $d,b,G,\varphi$ from
\begin{equation}
 \mathsf F_0=\left\{\gamma\in\R\;\middle|\;
 \begin{array}{l}
 \exists d\in\R^N\ \exists b\in\R\text{ such that }b<\gamma\text{ and}\\[2pt]
 \forall(G,\varphi)\in\Sym_\nu(\R)\times\R^{N+1},\\[2pt]
 \qquad\Psi_N^{\mathrm{sm}}(W(d,0),G,\varphi)\Rightarrow\varphi_N\le b
 \end{array}
 \right\}.
 \label{eq:GD-decision}
\end{equation}
The general ray identity~\eqref{eq:ray} gives $\mathsf F_0=(V_0,\infty)$, and Theorem~\ref{thm:smooth-memory} proves $V_0>c_N$ for $N\ge2$. In particular the threshold-set form of the strict memory gain is
\begin{equation}
 \mathsf F_1\setminus\mathsf F_0=(c_N,V_0]\ne\varnothing
 \qquad(N\ge2).
 \label{eq:smooth-threshold-gap}
\end{equation}
This is exactly the criterion of Corollary~\ref{cor:gain}. The argument computes the endpoints for every $m\ge1$ and proves every adjacent memory comparison. To obtain the exact algebraic value of $V_0$ for a specified $N\ge2$, one must still perform quantifier elimination on~\eqref{eq:GD-decision} and isolate its endpoint, as in Theorem~\ref{thm:procedure}. That computation is not carried out here, and no closed formula for $V_0$ as a function of $N$ is asserted.

\begin{table}[ht]
\centering
\renewcommand{\arraystretch}{1.45}
\begin{tabular}{@{}lccc@{}}
\toprule
Instance class & Horizon & $m_\star$ & $V_{\mathrm{lin}}$\\
\midrule
Convex quadratics & $N\ge1$ & $0$ & $\displaystyle\frac{LR^2}{2(2N+1)^2}$\\
Convex $L$-smooth & $N=1$ & $0$ & $\displaystyle\frac{LR^2}{8}$\\
Convex $L$-smooth & $N\ge2$ & $1$ & $\displaystyle\frac{LR^2}{2\theta_N^2}$\\
\bottomrule
\end{tabular}
\caption{Computed least auxiliary memory and unrestricted minimax value. Total persistent vector registers equal $m_\star+1$. The necessity of $m_\star=1$ in the last row uses the cited uniqueness theorem.}
\label{tab:results}
\end{table}

\FloatBarrier

\section*{Verification scope}
The accompanying Lean library is a partial formalization, pinned to Lean~4.35.0
and Mathlib~4.35.0. It verifies the realization equivalence in Theorem~\ref{thm:realization},
starting from the actual state recursion and formal-input correctness. The
construction and lower bound are proved for scalar inputs and extended to
one coefficient machine acting in every finite dimension. It also verifies
the cut/increment rank identity in
Proposition~\ref{prop:loci}, both inverse laws for the coefficient coordinates,
the equality $\Q_0=\{0\}$ on strictly lower triangular
matrices, and the threshold and gain statements in Theorem~\ref{thm:procedure}
and Corollary~\ref{cor:gain} for an explicitly supplied feasible-loss model and
memory set. The latter proofs use extended real suprema and infima, allow
infinite values, and do not assume attainment. The formal performance set is
constructed from loss inequalities, independently of the minimax value.
It additionally proves the projection graph and correction-readout identities,
query membership in $K$, singleton trajectories, and the Gram/PSD identities.
For quadratics it proves the residual-polynomial recurrence and its identity
with the matrix-gradient history trajectory, the exact spectral risk formula
with one-dimensional attainment, and the scaled Chebyshev candidate's degree,
normalization, and weighted interval maximum. The competing-polynomial lower
bound and the candidate's one-register implementation remain open in Lean.
For OGM it proves the actual two-register realization
and the exact least register count: one for $N=1$, two for $N\ge2$, when $L>0$.
It also proves that a finite minimax value is algebraic if a rational
quantifier-free formula is supplied and proved to describe its threshold set.
This last result does not construct the formula or perform elimination.
The invertible coordinate changes for register normalization and the identity
$F'_k\mathbf 1=\mathbf 1$ are proved, but the full trajectory conjugacy is not yet proved.
The library does not yet formalize normalized-register equivalence,
the general projected rank and trajectory equivalences, exact interpolation
models, quantifier elimination, or the benchmark optimality proofs. Those
claims are supported by the written proofs and cited results above.
An exact declaration map and a build/axiom audit accompany the source;
all audited declarations use only \texttt{propext}, \texttt{Classical.choice},
and \texttt{Quot.sound}.

The accompanying computations are supplementary finite checks \computation.
Exact rational arithmetic verifies the odd-Chebyshev polynomial certificate
for $N=1,\ldots,20$ and sampled cut-rank identities for $N=1,\ldots,8$.
Floating-point checks verify the OGM recurrence and coefficient calculations
for $N=1,\ldots,10$, and projected recurrences for $N=1,\ldots,6$, with tolerance
$10^{-9}$. These checks do not replace the all-horizon proofs, and the floating
point checks are not exact certificates. The programs, fixed inputs, and logs
are included in the accompanying source archive.

\section*{Provenance and AI disclosure}
\addcontentsline{toc}{section}{Provenance and AI disclosure}
All mathematical exposition, proofs, and computations were developed
during the Mathathon event window Nov~13, 8PM -- Nov~15, noon Pacific, 2026
\supportlabel{Event proof/computation}, except where explicitly cited as literature.
Coordination, editing, local proof development, and scripted checks used OpenAI Codex.
Within Open Math OS, coordination was configured for Kimi K3.1~Pro at
\texttt{reasoning\_effort=max}; the execution history also records coordinating
turns served by DeepSeek V4.1 Flash. Ultramath workers used
\texttt{meta/muse-spark-open} at \texttt{reasoning\_effort=max}. The authors directed
the work; independent reruns and audits are documented in the project, and the new
outputs await final human author review. No official Lean competition solution was used or copied.

\begin{thebibliography}{99}
\bibitem{BPR}
S.~Basu, R.~Pollack, and M.-F.~Roy.
\emph{Algorithms in Real Algebraic Geometry}.
Springer, 2003.
\href{https://doi.org/10.1007/978-3-662-05355-3}{doi:10.1007/978-3-662-05355-3}.

\bibitem{THG}
A.~B.~Taylor, J.~M.~Hendrickx, and F.~Glineur.
Smooth strongly convex interpolation and exact worst-case performance of first-order methods.
\emph{Mathematical Programming} \textbf{161} (2017), 307--345.
\href{https://arxiv.org/abs/1502.05666}{arXiv:1502.05666}.

\bibitem{KF}
D.~Kim and J.~A.~Fessler.
Optimized first-order methods for smooth convex minimization.
\emph{Mathematical Programming} \textbf{159} (2016), 81--107.
\href{https://arxiv.org/abs/1406.5468}{arXiv:1406.5468}.

\bibitem{Drori}
Y.~Drori.
The exact information-based complexity of smooth convex minimization.
\emph{Journal of Complexity} \textbf{39} (2017), 1--16.
\href{https://arxiv.org/abs/1606.01424}{arXiv:1606.01424}.

\bibitem{GJP}
B.~Grimmer, S.~Jo, and C.~Park.
On Minimax Optimality and Uniqueness of Fixed-Step First-Order Methods for Smooth Convex Optimization.
Preprint, September 2026, version~1.
\href{https://arxiv.org/abs/2609.07671v1}{arXiv:2609.07671v1}.

\bibitem{THGcomp}
A.~B.~Taylor, J.~M.~Hendrickx, and F.~Glineur.
Exact worst-case performance of first-order methods for composite convex optimization.
\emph{SIAM Journal on Optimization} \textbf{27} (2017), 1283--1313.
\href{https://doi.org/10.1137/16M108104X}{doi:10.1137/16M108104X}.

\bibitem{vdVeen95}
A.-J.~van der Veen.
Time-Varying Lossless Systems and the Inversion of Large Structured Matrices.
\emph{Archiv f\"ur Elektronik und \"Ubertragungstechnik} \textbf{49} (1995), 372--382.
Section~3.2, Theorem~2.
\href{https://sps.ewi.tudelft.nl/pubs/inverse95.pdf}{sps.ewi.tudelft.nl/pubs/inverse95.pdf}.

\bibitem{DewildevdVeen98}
P.~Dewilde and A.-J.~van der Veen.
\emph{Time-Varying Systems and Computations}.
Kluwer, 1998.
\href{https://doi.org/10.1007/978-1-4757-2817-0}{doi:10.1007/978-1-4757-2817-0}.
Book page: \url{https://sps.ewi.tudelft.nl/~allejan/kluwerbook.html}.

\bibitem{EidelmanGohberg05}
Y.~Eidelman and I.~Gohberg.
On generators of quasiseparable finite block matrices.
\emph{Calcolo} \textbf{42}(3--4) (2005), 187--214.
\href{https://doi.org/10.1007/s10092-005-0102-4}{doi:10.1007/s10092-005-0102-4}.

\bibitem{AST}
A.~d'Aspremont, D.~Scieur, and A.~B.~Taylor.
Acceleration Methods.
\emph{Foundations and Trends in Optimization} \textbf{5}(1--2) (2021), 1--245.
\href{https://doi.org/10.1561/2400000036}{doi:10.1561/2400000036}.

\bibitem{BnBPEP}
S.~Das Gupta, B.~P.~G.~Van Parys, and E.~K.~Ryu.
Branch-and-bound performance estimation programming: a unified methodology for constructing optimal optimization methods.
\emph{Mathematical Programming} \textbf{204} (2024), 567--639.
\href{https://doi.org/10.1007/s10107-023-01973-1}{doi:10.1007/s10107-023-01973-1}.

\bibitem{Vandenberghe22}
L.~Vandenberghe.
ECE236C --- Optimization Methods for Large-Scale Systems, Exercises.
UCLA, June~7, 2022, Exercise~17 (odd-Chebyshev polynomial for quadratic minimax).
\href{https://www.seas.ucla.edu/~vandenbe/236C/236C-exercises.pdf}{236C-exercises.pdf}.
\end{thebibliography}
\end{document}
